                         IMC 2016, Blagoevgrad, Bulgaria
                                       Day 1, July 27, 2016



Problem 1. Let f : [a, b] → R be continuous on [a, b] and differentiable on (a, b). Suppose that f
has infinitely many zeros, but there is no x ∈ (a, b) with f (x) = f 0 (x) = 0.
   (a) Prove that f (a)f (b) = 0.
   (b) Give an example of such a function on [0, 1].
                                       (Proposed by Alexandr Bolbot, Novosibirsk State University)

Solution. (a) Choose a convergent sequence zn of zeros and let c = lim zn ∈ [a, b]. By the continuity
of f we obtain f (c) = 0. We want to show that either c = a or c = b, so f (a) = 0 or f (b) = 0; then
the statement follows.
                                                                              f (zn ) − f (c)       0−0
    If c was an interior point then we would have f (c) = 0 and f 0 (c) = lim                 = lim        =
                                                                                  zn − c            zn − c
0 simultaneously, contradicting the conditions. Hence, c = a or c = b.
    (b) Let                                  
                                             x sin 1 if 0 < x ≤ 1
                                     f (x) =         x
                                             0          if x = 0.
                                          1
This function has zeros at the points       for k = 1, 2, . . ., and it is continuous at 0 as well.
                                         kπ
   In (0, 1) we have
                                                   1 1       1
                                         f 0 (x) = sin
                                                     − cos .
                                                   x x       x
Since sin x and cos x cannot vanish at the same point, we have either f (x) 6= 0 or f 0 (x) 6= 0 everywhere
           1        1

in (0, 1).

Problem 2. Let k and n be positive integers. A sequence (A1 , . . . , Ak ) of n × n real matrices is
preferred by Ivan the Confessor if A2i 6= 0 for 1 ≤ i ≤ k, but Ai Aj = 0 for 1 ≤ i, j ≤ k with i 6= j.
Show that k ≤ n in all preferred sequences, and give an example of a preferred sequence with k = n
for each n.
                                        (Proposed by Fedor Petrov, St. Petersburg State University)

Solution 1. For every i = 1, . . . , n, since Ai · Ai 6= 0, there is a column vi ∈ Rn in Ai such that
Ai vi 6= 0. We will show that the vectors v1 , . . . , vk are linearly independent; this immediately proves
k ≤ n.
    Suppose that a linear combination of v1 , . . . , vk vanishes:

                                 c1 v1 + . . . + ck vk = 0,    c1 , . . . , ck ∈ R.

For i 6= j we have Ai Aj = 0; in particular, Ai vj = 0. Now, for each i = 1, . . . , n, from
                                                              k
                                                              X
                          0 = Ai (c1 v1 + . . . + ck vk ) =         cj (Ai vj ) = ci (Ai vi )
                                                              j=1


we can see that ci = 0. Hence, c1 = . . . = ck = 0.
   The case k = n is possible: if Ai has a single 1 in the main diagonal at the ith position and its
other entries are zero then A2i = Ai and Ai Aj = 0 for i 6= j.
Remark. The solution above can be re-formulated using block matrices in the following way. Consider
                                               2                      
                                              A1       A1 0 . . . 0
                                                                2
                                           A2   0 A2 . . . 0 
                                                                       
                         A1 A2 . . . A k  .  =  .         .. . .    . .
                                             ..   ..       .     . .. 
                                                             Ak               0        0     . . . A2k
It is easy to see that the rank of the left-hand side is at most n; the rank of the right-hand side is at least k.
Solution 2. Let Ui and Ki be the image and the kernel of the matrix Ai (considered as a linear
operator on Rn ), respectively. For every pair i, j of indices, we have Uj ⊂ Ki if and only if i 6= j.
   Let X0 = Rn and let Xi = K1 ∩ K2 ∩ · · · ∩ Ki for i = 1, . . . , k, so X0 ⊃ X1 ⊃ . . . ⊃ Xk . Notice also
that Ui ⊂ Xi−1 because Ui ⊂ Kj for every j < i, and Ui 6⊂ Xi because Ui 6⊂ Ki . Hence, Xi 6= Xi−1 ;
Xi is a proper subspace of Xi−1 .
   Now, from
                              n = dim X0 > dim X1 > . . . > dim Xk ≥ 0
we get k ≥ n.

Problem 3. Let n be a positive integer. Also let a1 , a2 , . . . , an and b1 , b2 , . . . , bn be real numbers
such that ai + bi > 0 for i = 1, 2, . . . , n. Prove that
                                                       n         n
                                                                           n 2
                                                      P         P          P
                                 n
                               X ai b i − b       2       a i ·     b i −      bi
                                                  i   i=1       i=1        i=1
                                                    ≤            n                .
                                       a  i + b i
                                                                P
                                i=1                                (ai + bi )
                                                                        i=1

                  (Proposed by Daniel Strzelecki, Nicolaus Copernicus University in TorÞn, Poland)

Solution. By applying the identity
                                                  XY − Y 2       2Y 2
                                                           =Y −
                                                   X +Y         X +Y
                                                                                      n
                                                                                      P                   n
                                                                                                          P
with X = ai and Y = bi to the terms in the LHS and X =                                      ai and Y =          bi to the RHS,
                                                                                      i=1                 i=1

                            n                       n                                      n            n
                              ai b i − b 2                               2b2i                            b2i
                            X                       X                                      X            X
                                              i
                   LHS =                          =    bi −                           =     bi − 2           ,
                            i=1
                                  ai + b i            i=1
                                                                       ai + b i         i=1
                                                                                                      a + bi
                                                                                                   i=1 i

                                     n            n
                                                            n 2                                 n 2
                                     P            P         P                                     P
                                           ai ·       bi −    bi                 n
                                                                                 X                  bi
                                     i=1      i=1                i=1                           i=1
                        RHS =                 n
                                              P             n
                                                            P                =        bi − 2 P
                                                                                             n            .
                                                    ai +         bi               i=1          (ai + bi )
                                              i=1          i=1                                  i=1
Therefore, the statement is equivalent with
                                                                         n
                                                                        P    2
                                                  n                        bi
                                                  X       b2i         i=1
                                                                  ≥ P
                                                                    n            ,
                                                      a i +   b i
                                                  i=1                 (ai + bi )
                                                                       i=1

which is the same as the well-known variant of the Cauchy-Schwarz inequality,
                              n
                              X X2       i        (X1 + . . . + Xn )2
                                             ≥                                    (Y1 , . . . , Yn > 0)
                               i=1
                                      Yi            Y1 + . . . + Yn

with Xi = bi and Yi = ai + bi .
Problem 4. Let n ≥ k be positive integers, and let F be a family of finite sets with the following
properties:
   (i) F contains at least nk + 1 distinct sets containing exactly k elements;
                             

   (ii) for any two sets A, B ∈ F, their union A ∪ B also belongs to F.
   Prove that F contains at least three sets with at least n elements.
                                       (Proposed by Fedor Petrov, St. Petersburg State University)

Solution 1. If n = k then we have at least two distinct sets in the family with exactly n elements
and their union, so the statement is true. From now on we assume that n > k.
    Fix nk + 1 sets of size  k in F, call them ’generators’. Let V ∈ F be the union of the generators.
                           n
Since V has at least k + 1 subsets of size k, we have |V | > n.
    Call an element v ∈ V ‘appropriate’ if v belongs to at most n−1
                                                                                   
                                                                               k−1
                                                                                     generators. Then there exist
               n         n−1       n−1
                                     
at least k + 1 − k−1 = k + 1 generators not containing v. Their union contains at least n
elements, and the union does not contain v.
    Now we claim that among any n elements x1 , . . . , xn of V , there exists an appropriate element.
Consider all pairs (G, xi ) such that G is a generator       and xi ∈ G. Every generator has exactly k
elements, so the number of such pairs is at most nk + 1 · k. If some xi is not appropriate then xi
is contained in at least n−1       + 1 generators; if none of x1 , . . . , xn was appropriate then we wold have
                           k−1
at least n · k−1 + 1 pairs. But n · n−1
                   n−1
                                                   + 1 > ( n−1
                                                              
                                             k−1             k−1
                                                                   + 1) · k, so this is not possible; at least one
of x1 , . . . , xn must be appropriate.
    Since |V | > n, the set V contains some appropriate element v1 . Let U1 ∈ F be the union of all
generators not containing v1 . Then |U1 | ≥ n and v1 6∈ U1 . Now take an appropriate element v2 from
U1 and let U2 ∈ F be the union of all generators not containing v2 . Then |U2 | ≥ n, so we have three
sets, V , U1 and U2 in ∈ F with at least n elements: V 6= U1 because v1 ∈ V and v1 6∈ U1 , and U2 is
different from V and U1 because v2 ∈ V, U1 but v2 6∈ U2 .
Solution 2. We proceed by induction on k, so we can assume that the statement of the problem is
known for smaller values of k. By contradiction, assume that F has less than 3 sets with at least n
elements, that is the number of such sets is 0, 1 or 2. We can assume without loss of generality that
F consists of exactly N := nk + 1 distinct sets of size k and all their possible unions. Denote the
sets of size k by S1 , S2 , . . ..                                  S
    Consider a maximal set I ⊂ {1, . . . , N } such that A := i∈I Si has size less than n, |A| < n. This
means that adding any Sj for j ∈      / I makes the size at least n, |Sj ∪A| ≥ n. First,     let’s prove that such
                                                                                        |A|      n−1
                                                                                                     
j exist. Otherwise, all the sets Si are contained in A. But there are only k ≤ k < N distinct
k-element subsets of A, this is a contradiction. So there is at least one j such that |Sj ∪ A| ≥ n.
Consider all possible sets that can be obtained as Sj ∪ A for j ∈           / I. Their size is at least n, so their
number can be 1 or 2. If there are two of them, say B and C then B ⊂ C or C ⊂ B, for otherwise
the union of B and C would be different from both B and C, so we would have three sets B, C and
B ∪ C of size at least n. We see that in any case there must exist x ∈               / A such that x ∈ Sj for all
j∈ / I. Consider sets Sj0 = Sj \ {x} for j ∈  / I. Their sizes are equal to k − 1. Their number is at least
                                                                  
                                                n−1          n−1
                                         N−             =                + 1.
                                                   k         k−1
By the induction hypothesis, we can form 3 sets of size at least n − 1 by taking unions of the sets Sj0
for j ∈
      / I. Adding x back we see that the corresponding unions of the sets Sj will have sizes at least
n, so we are done proving the induction step.
    The above argument allows us to decrease k all the way to k =     0, so it remains to check the
statement for k = 0. The assumption says that we have at least n0 + 1 = 2 sets of size 0. This is
impossible, because there is only one empty set. Thus the statement trivially holds for k = 0.

Problem 5. Let Sn denote the set of permutations of the sequence (1, 2, . . . , n). For every permu-
tation π = (π1 , . . . , πn ) ∈ Sn , let inv(π) be the number of pairs 1 ≤ i < j ≤ n with πi > πj ; i.e. the
number of inversions in π. Denote by f (n) the number of permutations π ∈ Sn for which inv(π) is
divisible by n + 1.
                                                                          (p − 1)!
    Prove that there exist infinitely many primes p such that f (p − 1) >          , and infinitely many
                                                                             p
                                 (p − 1)!
primes p such that f (p − 1) <            .
                                     p
                                          (Proposed by Fedor Petrov, St. Petersburg State University)

Solution. We will use the well-known formula
                  X
                       xinv(π) = 1 · (1 + x) · (1 + x + x2 ) . . . (1 + x + · · · + xn−1 ).
                      π∈Sn

(This formula can be proved by induction on n. The cases n = 1, 2 are obvious. From each permu-
tation of (1, 2, . . . , n − 1), we can get a permutation of (1, 2, . . . , n) such that we insert the element
n at one of the n possible positions before, between or after the numbers 1, 2, . . . , n − 1; the number
of inversions increases by n − 1, n − 2, . . . , 1 or 0, respectively.)
    Now let                                             X
                                              Gn (x) =      xinv(π) .
                                                                       π∈Sn
                2πi
and let ε = e n+1 . The sum of coefficients of the powers divisible by n + 1 may be expressed as a
trigonometric sum as
                                            n−1                        n−1
                                        1 X                  n!    1 X
                          f (n) =               Gn (εk ) =      +          Gn (εk ).
                                      n + 1 k=0            n + 1 n + 1 k=1

Hence, we are interested in the sign of
                                                              n−1
                                                         n!   X
                                               f (n) −      =     Gn (εk )
                                                       n + 1 k=1

with n = p − 1 where p is a (large, odd) prime.
   For every fixed 1 ≤ k ≤ p − 1 we have
               p−1                                  p−1
           k
               Y
                         k   2k           (j−1)k
                                                    Y   1 − εjk   (1 − εk )(1 − ε2k ) · · · (1 − ε(p−1)k )
   Gp−1 (ε ) =     (1 + ε + ε + . . . + ε        )=           k
                                                                =                    k )p−1
                                                                                                           .
               j=1                                  j=1
                                                        1 − ε                 (1 − ε

Notice that the factors in the numerator are (1 − ε), (1 − ε2 ), . . . , (1 − εp−1 ); only their order is
different. So, by the identity (z − ε)(z − ε2 ) . . . (z − εp−1 ) = 1 + z + · · · + z p−1 ,
                                                    p                 p
                               Gp−1 (εk ) =           k p−1
                                                            =         2kπi p−1 .
                                            (1 − ε )            1−e p

Hence, f (p − 1) − (p−1)!
                     p
                          has the same sign as
                             p−1                                 p−1                             1−p
                             X              2kπi
                                                       1−p
                                                                 X         k(1−p)πi            πk
                                   (1 − e    p     )         =         e      p        −2i sin         =
                             k=1                                 k=1
                                                                                                p
                                                              p−1
                                                              2                                       1−p
                                      1−p               p−1   X       πk(p − 1)                   πk
                             =2·2           (−1)         2        cos                         sin             .
                                                              k=1
                                                                          p                        p
For large primes p the term with k = 1 increases exponentially faster than all other terms, so this
term determines the sign of the whole sum. Notice that cos π(p−1)    p
                                                                         converges to −1. So, the sum
is positive if p − 1 is odd and negative if p − 1 is even. Therefore, for sufficiently large primes,
f (p − 1) − (n−1)!
               p
                   is positive if p ≡ 3 (mod 4) and it is negative if p ≡ 1 (mod 4).
